\chapter{Gentle algebras, matching semigroups and determinant coordinates}
\label{ch:gentle-invariants}
\reviewstatus{The elementary decompositions, determinant transformations and worked semigroup presentation below have proofs in this chapter. General results on gentle semi-invariants and rational quotients are attributed to their primary sources with hypotheses retained. These are additional mathematical proofs, not additional Lean declarations.}

\section{From a relation to a variety of complexes}
Let \(k\) be an algebraically closed field of characteristic zero throughout this chapter, except where a result is explicitly stated over an arbitrary field. A bound quiver algebra \(A=kQ/I\) is \emph{gentle} when it is finite dimensional, every vertex has at most two incoming and two outgoing arrows, and the following restrictions hold. For any arrow, there is at most one possible successor whose composition is in \(I\), at most one successor whose composition is not in \(I\), and the analogous two restrictions hold for predecessors. The ideal \(I\) is generated by paths of length two. Thus two different kinds of combinatorics are present: paths that survive in \(A\), and paths whose consecutive pairs are relations. A drawing of a quiver without its relations does not determine the algebra.

The basic example is
\[
 1\xrightarrow{a}2\xrightarrow{b}3,\qquad ba=0.
\]
The algebra has basis \(e_1,e_2,e_3,a,b\) and dimension five. A representation of dimension vector \((p,q,r)\) is exactly a pair of matrices
\[
 k^p\xrightarrow{U}k^q\xrightarrow{V}k^r,\qquad VU=0.
\]
Consequently the representation scheme is the scheme of two consecutive differentials. Its defining equations, ranks, singular strata and coordinate ring are the objects studied in the preceding chapters. This is an equality of algebraic constructions, not an analogy between two diagrams.

\begin{proposition}[Complete module decomposition for the two-arrow relation]
Over any field, a representation \(k^p\xrightarrow{U}k^q\xrightarrow{V}k^r\) with \(VU=0\), \(\rank U=u\), and \(\rank V=v\), decomposes as
\[
 M_{12}^{\,u}\oplus M_{23}^{\,v}
 \oplus S_1^{\,p-u}\oplus S_2^{\,q-u-v}\oplus S_3^{\,r-v}.
 \tag{G1}\label{eq:gentle-a3-decomp}
\]
Here \(S_i\) is the simple module at vertex \(i\), while \(M_{12}\) and \(M_{23}\) have a one-dimensional space at the indicated two vertices and the identity on the joining arrow. These five modules are precisely the indecomposable modules.
\end{proposition}
\begin{proof}
Choose a complement \(P\) to \(\ker U\) in \(k^p\). Then \(U:P\to\im U\) is an isomorphism. Since \(\im U\subseteq\ker V\), choose
\[
 k^q=\im U\oplus H\oplus T,\qquad \ker V=\im U\oplus H.
\]
The restriction \(V:T\to\im V\) is an isomorphism. Complete \(\im V\) to a direct-sum decomposition of \(k^r\). Bases transported along the two isomorphisms exhibit the two collections of identity arrows. The remaining summands have all arrow maps zero and split into the displayed simples. Their multiplicities are the dimensions of the chosen spaces, giving (G1). Each of the five listed modules has endomorphism ring \(k\), so has no nontrivial idempotent and is indecomposable. The decomposition shows that no others occur.
\end{proof}

For dimension vector \((1,2,1)\), the four rank pairs give
\[
\begin{array}{c|l}
(u,v)&\text{module}\\ \hline
(0,0)&S_1\oplus S_2^{\,2}\oplus S_3\\
(1,0)&M_{12}\oplus S_2\oplus S_3\\
(0,1)&S_1\oplus S_2\oplus M_{23}\\
(1,1)&M_{12}\oplus M_{23}.
\end{array}
\]
Writing \(U=(x_1,x_2)^t\) and \(V=(y_1,y_2)\), the equation is \(x_1y_1+x_2y_2=0\). The rank-\((1,1)\) locus is dense in this three-dimensional quadric. In dimension \((1,1,1)\), however, the equation is \(xy=0\), and the two maximal rank pairs \((1,0)\), \((0,1)\) give two distinct irreducible components. Its coordinate ring \(k[x,y]/(xy)\) has zero divisors. It cannot be a subalgebra of a Laurent polynomial ring and therefore cannot be the ring of a cancellative affine semigroup in a lattice. This small example prevents an incorrect universal assertion that every complex scheme is itself an irreducible toric variety.

\subsection{Color chains and products}
Suppose now that \(Q\) has no oriented cycle. Starting at an arrow without a forbidden predecessor, follow the unique forbidden successor as long as one exists. The gentle restrictions and acyclicity ensure that this procedure terminates and partitions the arrows into directed chains. Color each such chain separately. The ideal \(I\) is generated by the monochromatic consecutive compositions.

For fixed vector spaces at the vertices, each color uses its own collection of matrix variables. Relations involving one color do not involve variables of another color. The polynomial algebra on all arrow matrices is the tensor product of the color polynomial algebras, and the relation ideal is the sum of the extended color ideals. Therefore
\[
 k[\operatorname{Rep}(A,\beta)]
   \cong\bigotimes_{\text{colors }s}
          k[\text{complex scheme along }s].
 \tag{G2}
\]
This assertion holds at the scheme level, including any scheme structure specified by the composition equations. It explains why straightening on complexes is integral to gentle invariant theory \citep[Sections 2--5]{smtCW2011}. However, two colors can meet at the same vertex, and the same \(\GL(\beta_x)\) acts on both. Taking invariants of (G2) therefore does not mean taking the tensor product of the separate invariant rings. The coupling of the two Schur factors at each shared vertex is the central additional step.

When oriented cycles are allowed, a forbidden chain may close up. It then yields a \emph{cyclic} complex, and the finite-chain statements must be reconsidered. Acyclicity in the quoted gentle-algebra theorems is a mathematical hypothesis, not an optional simplification of the picture.

\section{A second example: Kronecker pencils and regular modules}
The Kronecker quiver has two parallel arrows \(a,b:1\to2\), with no relations. Its path algebra is four dimensional because no two arrows compose. It is gentle. In dimension \((1,1)\), its representation space is \(k^2\), with coordinates \((a,b)\), and
\[
 (g_1,g_2)\cdot(a,b)=(g_2g_1^{-1}a,g_2g_1^{-1}b).
\]
All polynomial \(\GL_1\times\GL_1\)-invariants are constants: invariance under common scalar multiplication forces each nonconstant homogeneous part to vanish. The two coordinates are nevertheless semi-invariants with the same point-evaluation character \(g_2/g_1\). Since \(\SL_1\) is trivial, the total semi-invariant ring is \(k[a,b]\). Taking the projective quotient of its positive-degree part gives \(\mathbb P^1\). Thus an invariant ring that sees only constants and a semi-invariant ring that parametrizes a projective family are entirely compatible.

In dimension \((n,n)\), a representation is a matrix pencil \((A,B)\) with simultaneous left-right equivalence
\[
 (A,B)\longmapsto(g_2Ag_1^{-1},g_2Bg_1^{-1}).
\]
On the open locus \(\det A\ne0\), choose bases so that \(A=I_n\). The residual equivalence is conjugation of \(T=A^{-1}B\). Jordan decomposition consequently gives all indecomposable regular modules on this chart:
\[
 (I_n,J_n(\lambda)).
\]
Their endomorphism ring is \(k[t]/(t-\lambda)^n\), which is local, so they are indecomposable. The other chart includes \((J_n(0),I_n)\) at the point at infinity. For nonzero finite parameter, these are the usual band families; special zero and infinity presentations require the usual string/band conventions. The general string-algebra classification is due to Butler--Ringel and is not being reproved here \citep{GIButlerRingel1987}.

The polynomial
\[
 D_{A,B}(s,t)=\det(sA+tB)
\]
transforms by
\[
 D_{g_2Ag_1^{-1},g_2Bg_1^{-1}}(s,t)
       =\frac{\det g_2}{\det g_1}D_{A,B}(s,t).
 \tag{G3}\label{eq:pencil-weight}
\]
Every coefficient is a semi-invariant of the same character. Ratios of nonzero coefficients are rational invariants. They describe spectral parameters but cannot recover every extension. For example, \((I_2,\lambda I_2)\) is decomposable and \((I_2,J_2(\lambda))\) is indecomposable, yet both have determinant \((s+t\lambda)^2\). This is an explicit instance of the distinction between semisimplification data and the module itself. A proposed arithmetic determinant comparison must say which of these data it retains.

\section{Homological qualifications: a gentle algebra of infinite global dimension}
Consider the oriented triangle \(1\to2\to3\to1\), with every length-two path set to zero. This algebra is gentle and finite dimensional, with basis the three idempotents and three arrows. In the representation convention used here, the indecomposable projective \(P_i\) has top \(S_i\) and radical \(S_{i+1}\), with indices taken modulo three. Hence
\[
 \Omega S_i=S_{i+1},\qquad \Omega^3S_i=S_i.
 \tag{G4}
\]
Splicing the projective covers gives an infinite periodic resolution.

\begin{proposition}
For this algebra and \(n\ge0\),
\[
 \dim_k\Ext^n_A(S_i,S_j)
   =\begin{cases}1&j\equiv i+n\pmod3,\\0&\text{otherwise}.\end{cases}
\]
In particular the global dimension is infinite.
\end{proposition}
\begin{proof}
The \(n\)-th projective in the minimal resolution of \(S_i\) is \(P_{i+n}\). Applying \(\Hom_A(-,S_j)\) kills all differentials, since the resolution maps land in radicals and the target is semisimple. The term \(\Hom_A(P_{i+n},S_j)\) is \(k\) if the top is \(S_j\), and zero otherwise. These are therefore the cohomology groups. Nonzero groups in arbitrarily high degrees preclude finite global dimension.
\end{proof}
This example connects the periodic resolutions of orders to quiver algebra in a controlled way. It also shows why finite Euler sums and two-term determinant complexes require a projective-dimension hypothesis. One cannot use the hereditary formula \(\dim\Hom-\dim\Ext^1\) for every gentle module merely because its algebra is tame or its quiver is easy to draw.

\section{Schofield determinants, with the transformation law derived}
Let \(Q\) now be a finite acyclic quiver without imposed relations. Let \(V,W\) be representations of dimension vectors \(\alpha,\beta\). Define
\[
 C^0(V,W)=\bigoplus_x\Hom_k(V_x,W_x),\qquad
 C^1(V,W)=\bigoplus_{a:x\to y}\Hom_k(V_x,W_y)
\]
and
\[
 d_W^V(\phi)_a=W_a\phi_x-\phi_yV_a.
 \tag{G5}
\]
The kernel is \(\Hom_Q(V,W)\) by the definition of an intertwiner. The cokernel describes extensions.

\begin{proposition}[The elementary extension complex]
There is an exact sequence
\[
 0\longrightarrow\Hom_Q(V,W)\longrightarrow C^0(V,W)
 \xrightarrow{d_W^V} C^1(V,W)
 \longrightarrow\Ext^1_Q(V,W)\longrightarrow0.
 \tag{G6}\label{eq:schofield-exact}
\]
\end{proposition}
\begin{proof}
Any extension \(0\to W\to E\to V\to0\) splits as vector spaces at each vertex. Relative to these splittings its arrow maps have blocks
\[
 E_a=\begin{pmatrix}W_a&\eta_a\\0&V_a\end{pmatrix}.
\]
There are no path relations, so every collection \(\eta_a:V_x\to W_y\) defines a representation and hence an extension. Changing the vertex splittings by \(\phi_x:V_x\to W_x\) changes \(\eta\) by \(d_W^V(\phi)\), up to the fixed sign convention for the chosen change of basis. Conversely an equivalence of extensions fixing \(W\) and \(V\) has exactly this upper triangular form. Thus extension classes are the cokernel, and block addition realizes their vector-space addition. The kernel assertion was already observed.
\end{proof}

Set
\[
 \langle\alpha,\beta\rangle_Q
  =\sum_x\alpha_x\beta_x-\sum_{a:x\to y}\alpha_x\beta_y.
\]
If this is zero, (G5) is square and its determinant \(c^V(W)\) is defined after bases have been chosen. It is well defined up to a nonzero scalar. Equation (G6) proves
\[
 c^V(W)\ne0\quad\Longleftrightarrow\quad
 \Hom_Q(V,W)=0=\Ext^1_Q(V,W).
 \tag{G7}
\]
A zero of this determinant is therefore a concrete failure of exactness, rather than an unexplained spectral analogy.

Let \(g=(g_x)\in\prod_x\GL(W_x)\). On \(C^0\) let \(T_0\phi=(g_x\phi_x)\), and on \(C^1\) let \(T_1\eta=(g_y\eta_a)\). Direct substitution gives
\(d_{gW}^VT_0=T_1d_W^V\). Taking determinants proves
\[
 c^V(gW)=
 \prod_x\det(g_x)^{\,\sum_{a:y\to x}\alpha_y-\alpha_x}\ c^V(W).
 \tag{G8}\label{eq:schofield-weight}
\]
The exponent is minus the coefficient of \(\beta_x\) in the Euler form. Under the coordinate-function convention \(g\cdot f(W)=f(g^{-1}W)\), the weight changes sign, agreeing with the conventional weight \(\langle\alpha,\cdot\rangle_Q\). Stating the action removes an otherwise frequent sign ambiguity.

Derksen--Weyman's spanning theorem says that, for an acyclic quiver, these determinantal constructions span the semi-invariant weight spaces; their work also proves saturation results underlying Littlewood--Richardson multiplicities \citep{GIDerksenWeyman2000}. Those global theorems are external inputs. The exactness test and character computation needed for our examples have been proved directly.

\begin{example}[The Kronecker determinant]
Fix the representation \(V=(1,\lambda)\) of dimension \((1,1)\), and take \(W=(A,B)\) of dimension \((n,n)\). In compatible bases, (G5) is
\[
 d_W^V=\begin{pmatrix}A&-I_n\\B&-\lambda I_n\end{pmatrix},
 \qquad c^V(W)=\det(B-\lambda A).
 \tag{G9}
\]
For invertible \(A\), block elimination proves the determinant identity. Both sides are polynomials in the entries of \(A,B,\lambda\), and the invertible-\(A\) locus is dense, so the identity holds everywhere. For \(n=1\), the determinant \(b-\lambda a\) vanishes exactly when the two nonzero pairs have the corresponding regular parameter. For general \(n\), it detects a nonzero homomorphism from the chosen regular simple and hence an eigenvalue of the pencil, with the qualifications at infinity specified by the other chart.
\end{example}

For a bound quiver algebra, arrow blocks must also satisfy linearized relations, so the unrestricted \(C^1\) above no longer automatically computes extensions. A valid replacement is a projective resolution
\[
 0\longrightarrow P_1\xrightarrow{\delta}P_0\longrightarrow V\longrightarrow0
\]
when \(\operatorname{pd}V\le1\). Applying \(\Hom_A(-,W)\) then gives a two-term map whose kernel and cokernel are \(\Hom_A(V,W)\) and \(\Ext^1_A(V,W)\). If \(\delta\) is merely part of a noninjective presentation, its cokernel after applying \(\Hom\) need not be \(\Ext^1\): the next syzygy imposes another condition. The periodic triangle makes that warning testable.

\section{Schur factors and the semi-invariant matching condition}
The standard-monomial filtration assigns to each arrow \(a\) a partition \(\lambda(a)\), with length bounded by the relevant rank. At a vertex \(x\), each color \(s\) contributes a rational irreducible \(\GL(\beta_x)\)-representation of a dominant integral highest weight \(\gamma(x,s)\). Incoming partitions occur dualized and reversed; outgoing partitions occur with the opposite sign convention. The preceding chapter fixes the coordinate action and spells out these weights. A gentle vertex has at most two color factors, which is the decisive simplification.

\begin{proposition}[The two-factor character test]
Let \(L_\gamma,L_\nu\) be irreducible rational representations of \(\GL_d\) in characteristic zero. Then
\[
 (L_\gamma\otimes L_\nu)^{\SL_d}\ne0
 \quad\Longleftrightarrow\quad
 \gamma_i+\nu_{d+1-i}=c\quad(1\le i\le d)
 \tag{G10}
\]
for some integer \(c\). When nonzero this space is one dimensional, and \(\GL_d\) acts on it by \(\det^c\).
\end{proposition}
\begin{proof}
The invariant space identifies with
\[
 \Hom_{\SL_d}(L_\gamma^*,L_\nu).
\]
The dual has highest weight \((-\gamma_d,\ldots,-\gamma_1)\). Restriction to \(\SL_d\) is irreducible, and two rational \(\GL_d\)-highest weights restrict to the same \(\SL_d\)-highest weight exactly when their difference is an integer multiple of \((1,\ldots,1)\). Schur's lemma now gives the criterion and dimension. The one-dimensional \(\GL_d\)-subrepresentation has the stated determinant character.
\end{proof}
For a single color at a vertex, the condition reduces to a constant highest-weight tuple. Applying (G10) at every vertex defines an additive semigroup \(\Lambda_{\mathrm{SI}}\) of admissible arrow-partition data. The one-dimensional statement applies to each \emph{filtration factor}. Distinct partition data can have the same global character, so it does not assert that all semi-invariant weight spaces are one dimensional.

The strongest step is multiplicative. For an acyclic gentle algebra, a dimension vector \(\beta\), and a maximal admissible rank sequence \(r\) specifying an irreducible component, Carroll--Weyman construct normalized semi-invariants \(m_\lambda\) such that
\[
 m_\lambda m_\mu=m_{\lambda+\mu},\qquad
 k[\operatorname{Rep}(A,\beta)_r]^{\SL(\beta)}
       \cong k[\Lambda_{\mathrm{SI}}].
 \tag{G11}\label{eq:gentle-semigroup}
\]
This is their semi-invariant theorem \citep[Theorem 6.10]{smtCW2011}. It is stronger than saying that an associated graded vector space has Schur factors. Their construction uses the exterior diagonal maps splitting a top wedge into complementary wedges at a coupled vertex, then the maps that send wedge basis elements to minors. Straightening produces a distinguished standard multitableau with coefficient one, proving nonvanishing. A compatible column partition in the product yields the exact multiplicative normalization. The semigroup ring in (G11) is thus attached to actual semi-invariants on a specified component.

Subtracting adjacent entries in (G10) gives
\[
 \gamma_i-\gamma_{i+1}
       =\nu_{d-i}-\nu_{d+1-i}.
 \tag{G12}
\]
These nonnegative simple-root differences are propagated along a partition-equivalence graph. Its edges express equality of differences either at a shared vertex or along one arrow's incoming and outgoing copies. Components that are paths impose endpoint equations; components that are cycles leave a free nonnegative integer parameter. Writing \(u_a=\lambda(a)_{r(a)}\) for the last part, with zero-rank arrows omitted, the remaining semigroup has the form
\[
 U=\{u\in\N^{Q_1}:\varphi_e(u)=\varphi_{\Theta(e)}(u)\}.
 \tag{G13}
\]
Each endpoint expression is zero, one \(u_a\), or a sum \(u_a+u_b\); each variable occurs at most twice. The cycle parameters contribute polynomial variables:
\[
 k[\Lambda_{\mathrm{SI}}]\cong k[U][z_1,\ldots,z_b].
 \tag{G14}
\]
Here \(b\) counts cycles of the partition-equivalence graph. These combinatorial cycles should not be silently identified with band-module isomorphism classes or with handles of a Riemann surface; those objects have different definitions.

A matching graph depicts each equation as a dotted pairing of its two sides and each variable as a solid edge joining its occurrences, or as a solid loop if it occurs only once. Alternating walks account for balanced nonnegative exponent vectors. The reduced-equation condition means that a variable does not occur on \emph{both} sides of the same equation; it permits its coefficient to be zero on both sides. In particular it should not be read as requiring every pair of zero-one coefficients to be unequal. The matching-walk theorem yields generators with coordinate bounds two, and the gentle invariant degree bounds
\[
 \deg(\text{generators})\le 2\sum_{a\in Q_1}\binom{r(a)+1}{2},
 \qquad
 \deg(\text{relations})\le 8\sum_{a\in Q_1}\binom{r(a)+1}{2}.
 \tag{G15}
\]
These bounds refer to polynomial degree in arrow matrix entries, not just to the number of named generators in a binomial. General proofs are in \citet[Sections 7--9]{smtCW2011}. The following example supplies a fully proved instance.

\section{A matching semigroup with two explicit binomial relations}
Consider the semigroup \(S\subseteq\N^{10}\), with coordinates
\((a_1,\ldots,a_5,b_1,\ldots,b_5)\), defined by
\[
 \begin{aligned}
 a_1+a_2&=a_4+a_5,& b_1+b_2&=b_4+b_5,\\
 a_2+a_3&=b_2+b_3,& a_3+a_4&=b_3+b_4.
 \end{aligned}
 \tag{G16}
\]
Use the same letters for independent polynomial variables when writing its monomials. Define
\[
\begin{array}{ll}
 X_1=a_1a_5,&X_2=b_1b_5,\\
 Y_1=a_1a_4b_4b_1,&Y_2=a_5a_2b_2b_5,\\
 Z_1=a_2b_3a_4,&Z_2=b_2a_3b_4,\\
 B_1=a_2b_2b_4a_4,&B_2=a_3b_3.
\end{array}
 \tag{G17}
\]
These are the alternating-walk generators of the worked example in \citet[Example 8.20]{smtCW2011}. Direct substitution requires a correction to its arXiv-v1 printed relation: the factor multiplying \(X_1X_2\) is \(B_1\), not \(B_2\).

\begin{theorem}[Corrected explicit presentation]
The semigroup \(S\) is generated by the eight exponent vectors in (G17), and
\[
 k[S]\cong
 \frac{k[X_1,X_2,Y_1,Y_2,Z_1,Z_2,B_1,B_2]}
 {(Y_1Y_2-B_1X_1X_2,\ Z_1Z_2-B_1B_2)}.
 \tag{G18}\label{eq:matching-corrected}
\]
It is a normal affine domain of dimension six.
\end{theorem}
\begin{proof}
First take an arbitrary element of \(S\). The equations imply
\[
 d:=a_1-a_5=b_1-b_5=a_4-a_2=b_4-b_2.
\]
If \(d\ge0\), subtract \(d\) copies of the exponent vector of \(Y_1\); all four affected coordinates are at least \(d\). If \(d<0\), subtract \(-d\) copies of \(Y_2\). The residual coordinates satisfy \(a_1=a_5\), \(b_1=b_5\), \(a_2=a_4\), and \(b_2=b_4\). Remove the remaining \(a_1\) copies of \(X_1\) and \(b_1\) copies of \(X_2\).

Put \(A=a_2=a_4\), \(B=b_2=b_4\), \(C=a_3\), \(D=b_3\) in what remains. The single condition is \(A+C=B+D\). Subtract \(\min(A,B)\) copies of \(B_1\). If \(A\) remains positive, \(B=0\) and \(D=A+C\), so subtract \(A\) copies of \(Z_1\), followed by \(C\) copies of \(B_2\). If \(B\) remains positive, the symmetric procedure uses \(B\) copies of \(Z_2\) and then \(D\) copies of \(B_2\). If both are zero, use only \(B_2\). This proves generation with nonnegative coefficients.

The two displayed binomials vanish by multiplication of the ten underlying variables. To prove there are no additional defining relations, compute the integer kernel of the ten-by-eight exponent matrix. If the coefficients of a kernel vector in the displayed generator order are \(c_1,c_2,d_1,d_2,e_1,e_2,f_1,f_2\), the \(a_1,a_5,b_1,b_5\) coordinates give
\[
 d_1=d_2=p,\qquad c_1=c_2=-p.
\]
The \(a_3,b_3\) coordinates give \(e_1=e_2=q\), \(f_2=-q\), and the remaining coordinates give \(f_1=-p-q\). Thus every integer relation is
\[
 p(-1,-1,1,1,0,0,-1,0)
 +q(0,0,0,0,1,1,-1,-1).
\]
Reduce monomials using \(Y_1Y_2\mapsto B_1X_1X_2\) and
\(Z_1Z_2\mapsto B_1B_2\). The number of \(Y\)- and \(Z\)-factors strictly decreases, so reduction terminates. A reduced monomial contains neither both \(Y_1,Y_2\) nor both \(Z_1,Z_2\).
If two reduced monomials had the same underlying exponent vector, the kernel formula would apply. A nonzero \(p\) would force one of them to contain both \(Y\)'s after common factors are canceled, and a nonzero \(q\) would force both \(Z\)'s. Hence \(p=q=0\), and the reduced monomials are identical. Their images are therefore linearly independent. This proves the quotient presentation for arbitrary polynomials as well as for binomials.

The exponent matrix has rank six, since its integer kernel has rank two. The group algebra of its exponent lattice is a Laurent polynomial ring in six independent variables, so \(k[S]\) has dimension six. To check normality, write \(L=\operatorname{gp}(S)\subseteq\Z^{10}\). From (G16),
\[
 k[S]=k[L]\cap k[a_1,\ldots,a_5,b_1,\ldots,b_5]
\]
inside the ambient Laurent polynomial ring: an exponent in \(L\) with all coordinates nonnegative satisfies (G16), hence lies in \(S\). Both rings in the intersection are integrally closed. An element of the fraction field of \(k[S]\) integral over the intersection is integral over both rings and belongs to both fraction fields, hence belongs to the intersection. This proves normality. The monomial embedding proves that the ring is a domain.
\end{proof}

The first relation has degree eight in the underlying arrow-style variables, because \(\deg Y_i=4\), \(\deg B_1=4\), and \(\deg X_i=2\). The second has degree six, since \(\deg Z_i=3\) and \(\deg B_2=2\). Treating all eight symbols as degree one would destroy this grading and misstate the degree information. The companion verification checks every solution of (G16) whose ten coordinates lie in \(\{0,1,2\}\) by the constructive decomposition above, and separately verifies the two exponent identities. The finite check is additional evidence; the preceding proof establishes the unbounded assertion.

\section{Quotients and the information retained by semi-invariants}
Carroll--Chindris prove, for an acyclic gentle algebra over an algebraically closed field of characteristic zero, that the rational invariant field on each irreducible module-variety component is purely transcendental. Its transcendence degree is the sum of the multiplicities of the indecomposable regular components in the generic decomposition. For a regular irreducible component and a weight with nonempty semistable locus whose stable summands are regular, the corresponding moduli space is a product of projective spaces \citep[Theorem 1]{GICarrollChindris2013}. The qualifications about the component and stable summands are part of the theorem. They do not assert that every module variety, every affine invariant ring, or every stability quotient is a projective-space product.

Ratios of generalized Schofield determinants give their explicit rational coordinates. On the elementary Kronecker chart, this becomes familiar: the coefficients of \(\det(B-\lambda A)\), normalized by a nonzero coefficient, recover the symmetric functions of the generic regular parameters. Repeated eigenvalues and nonsplit extensions explain why these rational coordinates describe a generic quotient rather than a complete isomorphism classification at all points.

There are consequently three levels of geometry useful for this project. A representation variety remembers actual matrices satisfying relations. A semi-invariant ring remembers prescribed determinant characters and defines graded quotients. A rational invariant field describes a generic orbit quotient. The associated graded algebra of a straightening filtration adds another construction with its own specialization map. Passing between these objects is justified by maps and hypotheses, never by the fact that all four admit combinatorial descriptions.

\chapter{Exterior traces, zeta functions and the spectral comparison problem}
\label{ch:exterior-zeta}
\reviewstatus{The finite-dimensional determinant and trace identities are proved here. The cohomological description of elliptic-curve zeta functions and the Weil bound are external theorems. The transition to completed arithmetic \(L\)-functions is formulated as an obligation of the research program.}

\section{Exterior powers supply determinant coefficients}
Let \(T:V\to V\) be an endomorphism of an \(n\)-dimensional vector space over a characteristic-zero field. Its action on \(\bigwedge^jV\) is the special Schur functor associated to a column of height \(j\).

\begin{theorem}[Exterior determinant identity]
\[
 \det(I-tT)=\sum_{j=0}^n(-t)^j\Tr(\bigwedge\nolimits^j T).
 \tag{Z1}\label{eq:exterior-characteristic}
\]
\end{theorem}
\begin{proof}
For a diagonal matrix with eigenvalues \(\lambda_1,\ldots,\lambda_n\), the wedge basis has eigenvalues \(\lambda_{i_1}\cdots\lambda_{i_j}\). Its trace is therefore the \(j\)-th elementary symmetric polynomial, which is exactly the coefficient in \(\prod_i(1-t\lambda_i)\). For an arbitrary matrix, extend to an algebraic closure. The identity holds on matrices with distinct eigenvalues, a Zariski-dense open set. Each coefficient on both sides is a polynomial in the matrix entries, so the identity holds identically and descends to the original field.
\end{proof}
In particular the minors that generate standard monomials also compute the exterior traces in a characteristic polynomial. The scalar determinant is the top exterior character. General Schur functors organize more elaborate polynomial characters, but the existence of those characters by itself places no sign or location restriction on their roots.

For a \(T\)-invariant short exact sequence \(0\to U\to V\to W\to0\), a basis adapted to \(U\) makes \(T\) block upper triangular, proving
\[
 \det(I-tT_V)=\det(I-tT_U)\det(I-tT_W).
 \tag{Z2}
\]
No equivariant splitting is needed. In the Grothendieck group of finite-dimensional vector spaces equipped with an endomorphism, with relations from equivariant short exact sequences, this becomes multiplicativity on virtual differences:
\[
 D_{[V]-[W]}(t)=\frac{\det(I-tT_V)}{\det(I-tT_W)}.
 \tag{Z3}
\]
The denominator must be retained. An alternating determinant of a complex is naturally a rational function before additional cancellations are established.

\begin{proposition}[Cancellation through cohomology]
Let \((C^\bullet,d)\) be a finite complex of finite-dimensional vector spaces and let \(T\) be a chain endomorphism. Then
\[
 \prod_i\det(I-tT_{C^i})^{(-1)^i}
   =\prod_i\det(I-tT_{H^i})^{(-1)^i}.
 \tag{Z4}\label{eq:cohomological-determinant}
\]
\end{proposition}
\begin{proof}
Let \(Z^i=\ker d^i\) and \(B^i=\im d^{i-1}\). Both are \(T\)-stable. Apply (Z2) to
\(0\to Z^i\to C^i\to B^{i+1}\to0\) and
\(0\to B^i\to Z^i\to H^i\to0\).
Each boundary determinant appears once with exponent \(+1\) and once with exponent \(-1\) in the alternating product, leaving exactly the cohomology factors.
\end{proof}
This is the elementary determinant counterpart of the Euler characteristic identity. It states exactly what finite-complex exactness can do for a determinant: it cancels acyclic factors. It does not supply analytic continuation, an arithmetic Euler product or a positive inner product.

\section{Traces, closed walks and a finite zeta model}
In the formal power-series ring \(k[[t]]\), where characteristic zero allows division by positive integers,
\[
 -\log\det(I-tT)=\sum_{m\ge1}\frac{\Tr(T^m)}{m}t^m.
 \tag{Z5}
\]
To prove this, first diagonalize on the dense open set and use
\(-\log(1-t\lambda)=\sum_{m\ge1}(t\lambda)^m/m\). Each coefficient is a polynomial identity, so the same density argument proves it for all matrices. Alternatively differentiate the logarithm formally and expand \((I-tT)^{-1}\), then integrate with constant term zero. Over \(\C\), the same expression converges analytically for sufficiently small \(|t|\), for example when \(|t|\,\|T\|<1\).

For a finite directed graph with adjacency matrix \(B\), the entry \((B^m)_{vv}\) counts length-\(m\) walks from \(v\) to itself, and the trace counts closed walks with a marked starting point. Thus
\[
 Z_B(t):=\exp\left(\sum_{m\ge1}\frac{\Tr(B^m)}m t^m\right)
       =\det(I-tB)^{-1}.
 \tag{Z6}
\]
An Euler-product version follows by decomposing a closed walk into repetitions of a primitive cyclic word. A primitive cycle of length \(\ell\) has \(\ell\) marked starts, and its \(r\)-fold repetition contributes \(\ell/(r\ell)=1/r\), yielding the factor \((1-t^\ell)^{-1}\). These identities are formal: every coefficient receives finitely many contributions.

For a ribbon graph or dessin, choose \(B\) to be the nonbacktracking transition matrix on oriented edges if that is the walk model of interest. Its entries are one when two directed edges concatenate without immediately reversing, and zero otherwise. Formula (Z6) then defines the corresponding finite graph zeta function directly. The rotation system defining the Riemann surface and the transition rule defining \(B\) are separate inputs. Changing the former need not change the ordinary underlying graph, while changing the latter changes the trace counts. The interactive dessin therefore offers a precise combinatorial laboratory without identifying this finite rational function with a completed number-field \(L\)-function.

\section{The elliptic-curve numerator as a cohomological determinant}
Let \(E/\F_q\) be a smooth projective elliptic curve. With
\(N_m=\#E(\F_{q^m})\), define
\[
 Z(E,t)=\exp\left(\sum_{m\ge1}\frac{N_m}{m}t^m\right).
\]
The finite-field cohomological theorem and the Weil bound give
\[
 Z(E,t)=\frac{1-at+qt^2}{(1-t)(1-qt)},\qquad
 a=q+1-N_1,\quad |a|\le2\sqrt q.
 \tag{Z7}
\]
These are deep external facts; the elementary determinant identities above do not prove them. The endomorphism and Frobenius framework is described by Tate and the standard elliptic-curve sources already used in this monograph \citep{Tate1966Endomorphisms,SilvermanAEC,Sutherland2023Hasse}.

In a Frobenius convention compatible with \(N_m=1+q^m-\Tr(F^m|H^1)\), the numerator is
\(\det(I-tF|H^1)\), the two denominator factors come from \(H^0\) and \(H^2\), and (Z5) recovers the point-counting exponential. Equivalently, the roots \(\alpha,\beta\) of \(X^2-aX+q\) satisfy
\[
 \alpha+\beta=a,\quad \alpha\beta=q,\quad
 N_m=q^m+1-\alpha^m-\beta^m.
\]
The zeros in the \(t\)-plane are \(\alpha^{-1},\beta^{-1}\). Confusing eigenvalues with reciprocal zeros would reverse the radius.

For the companion curve \(y^2=x^3+2x+2\) over \(\F_{17}\), exact counting gives \(N_1=19\), so \(a=-1\). Its numerator is
\[
 1+t+17t^2,
\]
whose zeros are \((-1\pm i\sqrt{67})/34\) and have modulus \(1/\sqrt{17}\). The recurrence gives \(\alpha^2+\beta^2=a^2-2q=-33\), hence \(N_2=323\). The code independently counts the points over \(\F_{17^2}=\F_{17}[u]/(u^2-3)\) and verifies this value. This finite-field example directly connects the order generated by Frobenius, exterior traces, point counts, the ECC group and a zeta numerator. It does not turn the finite-field Weil theorem into the generalized Riemann hypothesis for number-field \(L\)-functions.

\section{What standard monomials can contribute to the arithmetic program}
The constructions now align in a useful sequence. A complex or quiver presentation specifies relations among maps. Exterior products provide minors, determinant characters and straightening generators. Schur filtrations keep track of the representation content of coordinate functions. At gentle vertices, complementary highest weights identify semi-invariants. Matching semigroups make a class of their rings and relations computable. Determinants of appropriate homomorphism complexes detect the failure of exactness. Exterior traces then turn an actual endomorphism into characteristic coefficients and a finite trace generating function.

Each step has a definite object and a proved or cited theorem. None yet identifies the result with an arithmetic \(L\)-function. To make that final identification, the inherited spectral program must still specify the arithmetic operator or complex, the exact Euler factors, local conductor and archimedean normalizations, the treatment of poles and zero multiplicities, and the analytic domain in which determinant identities are meaningful. In infinite dimension, an ordinary determinant usually does not exist; trace-class Fredholm determinants, zeta regularization and relative determinants have different hypotheses. A filtration identity or a formal power series at one point cannot silently supply equality after analytic continuation.

There is an elementary obstruction to any claim that gentle geometry alone forces the desired zero set. In the Kronecker example, take \(A=I_n\), \(B=\lambda I_n\) for an arbitrary complex \(\lambda\). All representation-theoretic relations remain valid, but
\(\det(B-zA)=(\lambda-z)^n\) has its root at any chosen \(\lambda\). If one imposes \(B=B^*\) in a specified positive Hermitian metric, the roots become real by the spectral theorem. If a generalized pencil is \(K-zM\) with \(M>0\) and \(K=K^*\), it reduces by \(M^{-1/2}\) to a self-adjoint eigenvalue problem. Positivity is additional structure on the representation, not a consequence of its Schur decomposition. The earlier positive-pencil implications state the corresponding conditional arithmetic conclusion precisely.

The same distinction controls the Yang--Mills comparison. The deformation complex and its determinant line describe zero modes, orientations and local moduli. A mass gap requires an inequality on the quantum Hamiltonian after the continuum and infinite-volume constructions. Standard monomial bases can simplify finite algebraic sectors and organize multiplicities, but they do not create a cutoff-uniform coercive estimate. The surface-code examples make the finite version visible: the binary chain condition proves stabilizer commutation, while code distance depends on the shortest nontrivial homology representatives, and a decoder's performance depends on a noise model and algorithm. These are compatible layers of one construction, with different proof obligations.

\section{A concrete program for further formalization}
The next formalization tasks can therefore be specified without suggesting that the cited theories are already encoded. One tractable target is the explicit monoid in (G16): define its ten integer coordinates, encode the constructive generator decomposition, and prove the two-binomial normal form. A second is the functorial determinant identity (Z4) for finite complexes and commuting chain maps. A third is the Schofield transformation law (G8), with the action convention in the type signature. These would extend the current verified chain-map/Pauli bridge with actual exterior and invariant-theory interfaces.

A much larger task is a reusable standard-monomial formalization: partitions and tableau orders, wedge shuffle maps over a stated base ring, a termination measure for straightening, standard-basis independence, the resulting filtration, and its equivariant Schur factors. In positive characteristic the distinction between a filtration and a semisimple decomposition must remain explicit. Formalizing only the rank identities or declaring the straightening theorem as a new axiom would not accomplish that task. The present repository supplies mathematical specifications and independently checked examples for it; its current verified Lean inventory is the 115 selected declarations compiled in the six-module package.
